\documentclass[12pt, a4paper]{article}

\usepackage[margin=1.8cm,top=2cm,bottom=2cm]{geometry}
\usepackage{amsmath,amssymb,amsfonts}
\usepackage{graphicx}
\usepackage[colorlinks=true,linkcolor=blue,citecolor=blue,urlcolor=blue]{hyperref}
\usepackage{bm,xcolor}
\usepackage[numbers,sort&compress]{natbib}
\usepackage{float}

\newcommand{\Geff}{G_{\mathrm{eff}}}
\newcommand{\cs}{c_s}
\newcommand{\meff}{m_{\mathrm{eff}}}
\newcommand{\heff}{\hbar_{\mathrm{eff}}}
\newcommand{\rhob}{\bar{\rho}}
\newcommand{\lP}{\ell_{\mathrm{P}}}
\newcommand{\mP}{m_{\mathrm{P}}}
\newcommand{\tP}{t_{\mathrm{P}}}

\begin{document}

\title{Singularity Regularization and the Emergent Planck Scale\\
in the Logarithmic Superfluid Vacuum}

\author{B. B. Kulangiev\\\small{Haifa, Israel}}
\date{April 2026}
\maketitle

\begin{abstract}
\noindent
General relativity predicts spacetime singularities where curvature invariants diverge. We show that the logarithmic superfluid vacuum framework provides a natural regularization mechanism through the quantum potential (Bohm potential) $Q = \Box\sqrt{\rho}/\sqrt{\rho}$, which was neglected at macroscopic scales in the companion paper but dominates at the soliton coherence scale. The healing length $\xi = \heff/(\meff c)$ of the condensate sets a minimum spatial resolution below which the hydrodynamic description breaks down and quantum pressure prevents gravitational collapse to a
point. We derive $\xi$ from the dimensional structure of the framework's gravitational coupling and show that
self-consistency---combining
$G \sim c^2/(\xi^2\rho_0)$ with the
healing-length identification
$\xi = \heff/(\meff c)$ and the natural
packing density $n_0 \sim 1/\xi^3$---uniquely
determines $\xi \sim \sqrt{\hbar G/c^3} = \lP$,
the Planck length, and identifies the condensate
particle mass as $\meff \sim \mP$ and the vacuum
mass density as $\rho_0 \sim \mP/\lP^3$, the
Planck density. The Planck scale thus emerges as the healing length of the vacuum condensate, not as a fundamental discreteness of spacetime. The temporal singularity is simultaneously resolved: the chemical potential $\mu(x) = \meff c_s^2/\heff$, which governs the local clock rate, is bounded by the same mechanism,
ensuring that gravitational time dilation is always
finite and the phase of the condensate remains
well-defined everywhere. The corresponding cutoffs
$\Delta x_{\min} \sim \lP$ and
$\Delta t_{\min} \sim \tP$ are related by $c$,
respecting the emergent Lorentz invariance they
regularize. These identifications unify $\hbar$,
$c$, and $G$ as properties of a single physical
substrate: the speed of sound, the quantum coherence
scale, and the emergent coupling of the vacuum
condensate.
\end{abstract}


% ====================================================================
\section{Introduction}
\label{sec:intro}
% ====================================================================

The Penrose-Hawking singularity theorems~\cite{penrose_1965,hawking_1970}
establish that geodesic incompleteness is a generic
feature of general relativity under physically
reasonable energy conditions. At a singularity, the
curvature invariants diverge, the metric ceases to be
defined, and the predictive power of the theory breaks
down. This is widely regarded as a signal that GR must
be replaced by a more fundamental theory at sufficiently
small scales~\cite{dewitt_1967}.

In quantum gravity approaches, the Planck length
$\lP = \sqrt{\hbar G/c^3} \approx 1.616 \times 10^{-35}$~m
is typically introduced as a fundamental
discreteness scale or minimum length, either through
dimensional analysis or through specific mechanisms
such as loop quantization~\cite{rovelli_2004, ashtekar_2006} or string theory~\cite{polchinski_1998}. However, the
physical origin of $\lP$ as a minimum length
remains unclear: why should spacetime have a minimum
resolution, and why at precisely this scale~\cite{garay_1995}?

In a companion paper~\cite{kulangiev_2026a}, we developed a superfluid vacuum model based on the relativistic logarithmic Klein-Gordon equation, deriving emergent gravitational light bending with PPN parameter $\gamma = 1$. The framework treats the vacuum as a complex order parameter $\psi(x,t)$ whose Madelung decomposition $\psi = \sqrt{\rho}\,e^{iS/\heff}$ produces hydrodynamic equations with a density-dependent sound speed and a macroscopic flow velocity. Throughout that paper, the quantum potential (Bohm potential)~\cite{bohm_1952}
$Q = \Box\sqrt{\rho}/\sqrt{\rho}$ was neglected
in the macroscopic limit, valid at length scales much
larger than the condensate's coherence scale.

In this paper, we restore $Q$ and show that it
provides a natural singularity regularization
mechanism. The healing length $\xi = \heff/(\meff c)$
of the condensate---the scale at which $Q$ balances
the logarithmic interaction---sets a minimum spatial
resolution. We then show that self-consistency of the
framework, requiring the dimensionally-correct gravitational
coupling $G \sim c^2/(\xi^2 \rho_0)$ to match the
observed Newton's constant when combined with the
healing-length structure of the soliton and the natural
condensate packing density $n_0 \sim 1/\xi^3$, uniquely
fixes $\xi = \lP$, the Planck length.


% ====================================================================
\section{The Quantum Potential at Short Distances}
\label{sec:bohm}
% ====================================================================

The Madelung decomposition~\cite{madelung_1927} of the logarithmic Klein-Gordon equation~\cite{kulangiev_2026a} yields the Hamilton-Jacobi equation (in energy units, so that both sides have dimensions of energy):
\begin{equation}
    \frac{1}{2\meff}\,\partial_\mu S\,\partial^\mu S
    = \meff c^2 - b\,\ln\!\left(\frac{\rho}{\rho_0}\right)
    + Q[\rho]\,,
\label{eq:hj}
\end{equation}
where $b$ is the logarithmic coupling with units of
energy, and the quantum potential (Bohm potential) is
\begin{equation}
    Q[\rho] = -\frac{\heff^2}{2\meff}\,\frac{\Box\sqrt{\rho}}{\sqrt{\rho}}\,.
\label{eq:bohm}
\end{equation}
$Q$ has units of energy, consistent with the other terms.
In the companion paper, $Q$ was neglected at macroscopic
scales ($r \gg \xi$), where the density varies slowly.
We now examine the regime $r \sim \xi$ where $Q$
becomes dominant.

\subsection{Scaling analysis}

For a density profile varying on a length scale $L$,
the quantum potential (Eq.~\ref{eq:bohm}) scales as:
\begin{equation}
    Q \sim \frac{\heff^2}{\meff\,L^2}\,.
\label{eq:Q_scaling}
\end{equation}
The gravitational interaction term (from the logarithmic nonlinearity) scales as:
\begin{equation}
    b\,\ln\!\left(\frac{\rho}{\rho_0}\right)
    \sim b\,\frac{\delta\rho}{\rho_0}\,,
\label{eq:grav_scaling}
\end{equation}
which for a weak-field gravitating source becomes
$b\,\delta\rho/\rho_0 \sim b\,GM/(c^2 r)$, using the
Bernoulli relation~\cite{kulangiev_2026a}.

The quantum potential dominates over the logarithmic
interaction when $Q \sim b$, i.e., when
$\heff^2/(\meff\,L^2) \sim b$. This defines a crossover
length scale:
\begin{equation}
    L_* = \sqrt{\frac{\heff^2}{\meff\,b}}\,.
\end{equation}
Identifying $b = \meff c^2$ at the relativistic
background (so that the logarithmic term has magnitude
comparable to the rest energy), we obtain the healing
length:
\begin{equation}
    \boxed{\;\xi \equiv \frac{\heff}{\meff\,c}\;}
\label{eq:healing}
\end{equation}
consistent with~\cite{kulangiev_2026a}. At
$r \lesssim \xi$, the quantum pressure from $Q$
overwhelms the interaction term, preventing the density
from diverging.

\subsection{The repulsive core}

The physical mechanism is identical to that which
stabilizes vortex cores in laboratory
superfluids~\cite{volovik_2003}: the quantum pressure
creates a repulsive force that increases faster than
the attractive force as the core is compressed.
Quantitatively, for a radially symmetric density
profile $\rho(r)$:
\begin{equation}
    Q = \frac{1}{\sqrt{\rho}}\left(
    \frac{d^2\sqrt{\rho}}{dr^2}
    + \frac{2}{r}\frac{d\sqrt{\rho}}{dr}\right)\,.
\end{equation}
If $\rho \to \rho_{\max}$ as $r \to 0$ (a regularized
core), then $Q \to Q_0 = \text{const}$ at the center.
If instead $\rho \to \infty$ (a singularity), then
$Q \to -\infty$, providing an arbitrarily strong
repulsive force. The singularity is self-regulating:
any attempt to compress the density beyond $\rho_{\max}$
generates a quantum pressure that exceeds the
gravitational compression.


% ====================================================================
\section{The Gausson Ground State}
\label{sec:gausson}
% ====================================================================

The ground-state soliton of the logarithmic wave
equation is the Gausson---a Gaussian-localized
solution discovered by Bia{\l}ynicki-Birula and
Mycielski~\cite{bialynicki_birula_1979}. For the
three-dimensional stationary equation with the
logarithmic potential, the density profile is:
\begin{equation}
    \rho(r) = \rho_{\max}\,
    \exp\!\left(-\frac{r^2}{2\sigma^2}\right)\,,
\label{eq:gausson}
\end{equation}
where $\sigma$ is the Gausson width. Substituting
into the stationary Hamilton-Jacobi equation with
$Q$ retained yields the width:
\begin{equation}
    \sigma = \frac{\heff}{\sqrt{2\,\meff\,b}}\,.
\label{eq:sigma}
\end{equation}
With the relativistic identification $b \sim \meff c^2$
from the balance between the logarithmic potential and
the rest energy in Eq.~(\ref{eq:hj}):
\begin{equation}
    \sigma \sim \frac{\heff}{\meff\,c} = \xi\,.
\label{eq:sigma_xi}
\end{equation}
The Gausson width equals the healing length: the
soliton core size IS the minimum spatial resolution
of the condensate.

The quantum potential at the Gausson center evaluates to:
\begin{equation}
    Q(0) = \frac{3\,\heff^2}{2\,\meff\,\sigma^2}
    = \frac{3\,\meff\,c^2}{2}\,,
\label{eq:Q_center}
\end{equation}
which is of order the rest energy $\meff c^2$. The
associated curvature scale is $Q(0)/\heff^2 \sim 1/\xi^2$. This is the maximum curvature the acoustic metric can sustain: the Kretschmann scalar~\cite{kretschmann_1917} of the acoustic metric
is bounded by $K \lesssim 1/\xi^4$.

% ====================================================================
\newpage
\section{Derivation of the Healing Length from Known Parameters}
\label{sec:derivation}
% ====================================================================

We now show that $\xi$ is determined by the parameters
derived in this paper from dimensional analysis. The chain of constraints is:

\subsection{Dimensional structure of the gravitational coupling}

For an emergent gravity sourced by perturbations of a
condensate of mass density $\rho_0$, dimensional analysis fixes the form of Newton's constant. The available parameters are the sound speed $c$, the condensate mass density $\rho_0$, and a length scale $\xi$ characterizing the response of the condensate to a localized perturbation. The unique combination with the dimensions
of $G = [\mathrm{m}^3/(\mathrm{kg\,s}^2)]$ is:
\begin{equation}
    G \;\sim\; \frac{c^2}{\xi^2\,\rho_0}\,,
\label{eq:G}
\end{equation}
up to a dimensionless prefactor of order unity. The
length scale $\xi$ in the denominator must enter
squared by inversion of the Laplacian in the linearized
condensate response; we will identify this with the
healing length below.

\subsection{The Gausson energy}

The Gausson rest energy is given by the
Bia{\l}ynicki-Birula--Mycielski energy
functional~\cite{bialynicki_birula_1979} of the
logarithmic Schr\"{o}dinger equation. For a Gaussian
density profile $\rho(r) = n_0\,e^{-r^2/(2\sigma^2)}$
of width $\sigma = \xi = \heff/(\meff c)$, the
quantum-pressure contribution to the energy is, at
the order of magnitude,
\begin{equation}
    E_v \sim \int Q\,\rho\,d^3x
    \sim \frac{\heff^2}{\meff\,\xi^2}\cdot\xi^3\,n_0
    = \frac{\heff^2\,n_0\,\xi}{\meff}\,,
\label{eq:Ev}
\end{equation}
where $n_0$ is the central condensate number density
(units $\mathrm{m}^{-3}$) and $Q\sim\heff^2/(\meff\,\xi^2)$
is the quantum potential evaluated on the
Gausson~(Eq.~\ref{eq:Q_center}). Substituting
$\xi = \heff/(\meff c)$:
\begin{equation}
    E_v \sim \frac{\heff^3\,n_0}{\meff^2\,c}\,.
\label{eq:Ev2}
\end{equation}
Dimensional check: $[\heff^3 n_0/(\meff^2 c)] =
(\mathrm{J\cdot s})^3 \cdot \mathrm{m}^{-3} /
(\mathrm{kg}^2 \cdot \mathrm{m/s}) = \mathrm{J}$ \checkmark.

\subsection{Self-consistency}
\label{sec:self_consistency}

We combine three relations:
\begin{itemize}
\item the gravitational coupling
  Eq.~(\ref{eq:G}): $G \sim c^2/(\xi^2 \rho_0)$\,,
\item the healing-length structure of the soliton
  Eq.~(\ref{eq:healing}): $\xi = \heff/(\meff c)$\,,
\item the relation between mass density and number
  density: $\rho_0 = \meff\,n_0$\,.
\end{itemize}
Substituting $\xi$ and $\rho_0$ into Eq.~(\ref{eq:G}):
\begin{equation}
    G \;\sim\;
    \frac{c^2}{\xi^2\,\meff\,n_0}
    \;=\;\frac{c^2 \cdot \meff^2 c^2}{\heff^2\,\meff\,n_0}
    \;=\;\frac{\meff\,c^4}{\heff^2\,n_0}\,,
\label{eq:G_chain}
\end{equation}
which can be solved for $\meff$ in terms of $n_0$:
\begin{equation}
    \meff \;\sim\; \frac{G\,\heff^2\,n_0}{c^4}\,.
\label{eq:meff}
\end{equation}
Dimensional check: $[G\heff^2 n_0/c^4] =
(\mathrm{m}^3/(\mathrm{kg\,s}^2))(\mathrm{J\,s})^2(\mathrm{m}^{-3})/(\mathrm{m/s})^4
= \mathrm{kg}$ \checkmark.

\subsection{Closure: \texorpdfstring{$\heff = \hbar$ and $n_0 \sim 1/\xi^3$}{heff = hbar and n0 ~ 1/xi cubed}}

Equation~(\ref{eq:meff}) contains three unknowns:
$\heff$, $\meff$, and $n_0$. We fix them by two
natural identifications.

First, $\heff = \hbar$: the condensate's quantum
mechanics reduces to standard quantum mechanics for
its excitations (phonons identified with photons,
solitons with particles).

Second, $n_0 \sim 1/\xi^3$: the condensate has one
particle per healing volume. This is the natural
packing density for a self-interacting relativistic
condensate whose healing length equals the minimum
spatial resolution. With $\xi = \hbar/(\meff c)$:
\begin{equation}
    n_0 \;\sim\; \frac{1}{\xi^3}
    \;=\;\frac{\meff^3 c^3}{\hbar^3}\,.
\label{eq:n0}
\end{equation}
Substituting~(\ref{eq:n0}) into~(\ref{eq:meff}):
\begin{equation}
    \meff \;\sim\; \frac{G\,\hbar^2}{c^4}
    \cdot \frac{\meff^3 c^3}{\hbar^3}
    \;=\;\frac{G\,\meff^3}{\hbar\,c}\,,
\end{equation}
which simplifies to:
\begin{equation}
    \meff^2 \;\sim\; \frac{\hbar\,c}{G}\,.
\label{eq:meff_closure}
\end{equation}
The right-hand side is precisely $\mP^2$, where
$\mP = \sqrt{\hbar c/G}$ is the Planck mass.
Therefore:
\begin{equation}
    \boxed{\;\meff \sim \mP\;,\qquad
    \xi \sim \lP\;,\qquad
    n_0 \sim 1/\lP^3\;}
\label{eq:planck_identifications}
\end{equation}
The condensate particle mass is the Planck mass, the
healing length is the Planck length, and the vacuum
number density is one particle per Planck volume.
The $O(1)$ prefactors depend on the detailed Gausson
energy integral; the essential result is that all
three scales are fixed at the Planck scale by
self-consistency.

\subsection{The healing length IS the Planck length}

Substituting the Planck-mass identification~(\ref{eq:planck_identifications})
into $\xi = \hbar/(\meff c)$:
\begin{equation}
    \boxed{\;\xi\sim\frac{\hbar}{\mP\,c}=\lP\;}
\label{eq:xi_planck}
\end{equation}
where $\lP = \sqrt{\hbar G/c^3} \approx 1.616 \times
10^{-35}$~m. The Planck length emerges as the healing
length of the vacuum condensate, determined by
self-consistency of the framework---not assumed. An
equivalent statement is that the vacuum condensate has
one particle per Planck volume, giving a Planck-density
reservoir.

% ====================================================================
\section{The Spacetime UV Cutoff}
\label{sec:uv}
% ====================================================================

\subsection{Minimum spatial resolution}

Below the healing length $\xi \approx \lP$, the
hydrodynamic (Madelung) description breaks down:
the quantum potential $Q$ dominates all other terms,
and the system is governed by the full wave equation
rather than the emergent acoustic metric. The acoustic
metric---and with it, the emergent spacetime
geometry---ceases to be well-defined at scales
$\Delta x \lesssim \xi$.

This is not a discretization of spacetime. The
condensate wave function $\psi(x,t)$ remains
continuous and well-defined at all scales. What has
a minimum resolution is the \emph{emergent} spacetime
geometry experienced by collective excitations
(phonons/photons). The underlying substrate is
continuous; the emergent geometry is coarse-grained.

\subsection{Minimum temporal resolution}

The phase of the condensate oscillates as
$S(t) = -\meff c^2 t/\heff$ in the rest
frame~\cite{kulangiev_2026a}. The minimum temporal
interval resolvable by the acoustic metric is the
time for a phase oscillation to cross the healing
length:
\begin{equation}
    \Delta t_{\min} = \frac{\xi}{c}
    = \frac{\hbar}{\meff\,c^2}
    \sim \frac{\lP}{c} = \tP\,,
\label{eq:dt_min}
\end{equation}
where $\tP = \sqrt{\hbar G/c^5} \approx 5.39 \times
10^{-44}$~s is the Planck time.

\subsection{Unified UV cutoff}

The spatial and temporal cutoffs are related by the
sound speed:
\begin{equation}
    \frac{\Delta x_{\min}}{\Delta t_{\min}} = c\,.
\label{eq:uv_ratio}
\end{equation}
This is a consequence of the Lorentz invariance of
the acoustic metric at scales $\gg \xi$. The UV
cutoff respects the emergent symmetry that it
regularizes.

\subsection{Maximum curvature}

The Kretschmann scalar of the acoustic metric is
bounded:
\begin{equation}
    K = R_{\mu\nu\alpha\beta}\,R^{\mu\nu\alpha\beta}
    \lesssim \frac{1}{\xi^4} \approx \frac{1}{\lP^4}\,.
\label{eq:Kretschmann}
\end{equation}
This is the same bound obtained in loop quantum
gravity~\cite{rovelli_2004}, but here it arises from
the condensate microphysics rather than from a
quantization of geometry.


% ====================================================================
\section{Temporal Regularization: The Phase Counter}
\label{sec:temporal}
% ====================================================================

\subsection{Time as phase accumulation}

In the Madelung representation, the condensate phase
$S(x,t)$ satisfies the Hamilton-Jacobi
equation~(\ref{eq:hj}). In the rest frame of the
condensate, $S(t) = -\mu\,t$, where the chemical
potential
\begin{equation}
    \mu(x) = \frac{\meff\,c_s^2(x)}{\heff}
\label{eq:mu}
\end{equation}
sets the local oscillation rate of the order parameter
$\psi \propto e^{-i\mu t}$. This phase accumulation IS
the local clock: two events separated by $\Delta t$
accumulate a phase difference $\Delta\phi = \mu\,\Delta t$.
A clock measures time by counting phase cycles of a
physical oscillator; here, the oscillator is the vacuum
itself.

Gravitational time dilation is the spatial variation
of $\mu(x)$. In the weak-field limit, the acoustic
metric gives $g_{00} \propto -c_s^2(x) \propto
-\rho(x)$, so:
\begin{equation}
    \frac{d\tau}{dt} = \sqrt{-g_{00}}
    \propto c_s(x) \propto \sqrt{\rho(x)}\,.
\label{eq:time_dilation}
\end{equation}
Deeper in a gravitational well (lower $\rho$), the
sound speed is lower, the phase accumulates more
slowly, and clocks run slower---exactly the
gravitational redshift.

\subsection{The temporal singularity in GR}

In general relativity, the singularity at $r = 0$
inside a black hole is not merely a point of
infinite density; it is a boundary where time itself
ceases to exist. The Kretschmann scalar diverges,
the metric is undefined, and no further time
evolution is possible. At the event horizon ($r = r_H$),
an external observer sees clocks approach infinite
redshift: $g_{00} \to 0$, so $d\tau/dt \to 0$.

Both pathologies involve the metric component $g_{00}$
reaching extreme values: zero at the horizon
(for external coordinates) and undefined at the
singularity.

\subsection{Regularization by the bounded chemical potential}

In the superfluid framework, $g_{00} \propto
-\rho\,c_s \propto -\rho^{3/2}$. Since the Bohm
potential bounds the density:
\begin{equation}
    0 < \rho_{\min} \leq \rho(x)
    \leq \rho_{\max} \sim e\,\rho_0\,,
\label{eq:rho_bounds}
\end{equation}
the chemical potential is bounded from both sides:
\begin{equation}
    0 < \mu_{\min} \leq \mu(x) \leq \mu_{\max}\,.
\label{eq:mu_bounds}
\end{equation}
The clock never stops ($\mu > 0$ everywhere) and
never runs infinitely fast ($\mu < \mu_{\max}$).

The period of the fastest possible phase
oscillation of the condensate is set by the maximum
chemical potential:
\begin{equation}
    T_{\mathrm{phase}} = \frac{2\pi}{\mu_{\max}}
    = \frac{2\pi\,\heff}{\meff\,c^2}
    \;\sim\; 2\pi\,\tP\,.
\label{eq:dt_phase}
\end{equation}
This is consistent with the resolution limit
$\Delta t_{\min} \sim \tP$ derived from the
sound-crossing time of the healing length
(Eq.~\ref{eq:dt_min}); the two definitions differ
by a factor of $2\pi$ corresponding to one full
phase cycle versus a half-period crossing. No
physical process within the emergent spacetime can
resolve time intervals shorter than $\sim \tP$.

The maximum resolvable time dilation ratio
(the gravitational redshift between the deepest
well and the cosmological background) is:
\begin{equation}
    \frac{\mu_{\max}}{\mu_{\min}}
    = \frac{c_{s,\max}^2}{c_{s,\min}^2}
    = \frac{\rho_{\max}}{\rho_{\min}}
    \leq \frac{e\,\rho_0}{\rho_{\min}}\,.
\label{eq:max_redshift}
\end{equation}
This is finite. There is no infinite redshift
surface---the horizon is replaced by a region of
large but finite time dilation, across which the
condensate phase remains well-defined and
continuous.

\subsection{The arrow of time}

The phase of the condensate advances monotonically:
$S(t+\Delta t) = S(t) - \mu\,\Delta t$ with
$\mu > 0$. This provides a microscopic arrow of
time that is absent in the classical wave equation
(which is time-reversal symmetric). The phase winds
monotonically in one direction because $\mu$ is
positive-definite throughout the condensate, and
this directional phase accumulation defines the
forward direction of time.

This connects the thermodynamic arrow (entropy
increases toward equilibrium with the Bulk) to
the cosmological arrow (the Universe expands as
the condensate evolves) and to the quantum-mechanical
arrow (the phase accumulates monotonically).

% ====================================================================
\newpage
\section{Singularity Regularization}
\label{sec:singularity}
% ====================================================================

\subsection{Black holes as phase boundaries}

In the acoustic metric framework, the event horizon
occurs where the inflow velocity equals the sound
speed: $v(r_H) = c$. The Schwarzschild radius is:
\begin{equation}
    r_H = \frac{2GM}{c^2}\,.
\end{equation}
For a solar-mass object, $r_H \approx 3$~km $\gg \xi$.
The acoustic metric description is valid at and
outside the horizon.

Inside the horizon, the inflow velocity exceeds $c$,
and the acoustic metric describes a trapped region.
However, as $r \to 0$, the density profile of the
infalling condensate steepens until the gradient
reaches the healing-length scale. At $r \sim \xi$,
the quantum potential~(\ref{eq:Q_center}) becomes
comparable to the gravitational term, and the
density is regularized to a finite maximum:
\begin{equation}
    \rho_{\max} \sim \rho_0\,
    \exp\!\left(\frac{\meff c^2}{\heff^2/(\meff\xi^2)}\right)
    = \rho_0\,e^{\,\meff^2 c^2\xi^2/\heff^2}
    = \rho_0\,e\,,
\label{eq:rho_max}
\end{equation}
where the last equality uses $\xi = \heff/(\meff c)$
so that the exponent equals $1$. The singularity is replaced by a Planck-scale core of radius $\sim \xi \sim \lP$ operating at the maximum vacuum density $e\,\rho_0$, where the hydrodynamic description gives way to the full quantum dynamics of the condensate.

\subsection{Comparison with other approaches}

The regularization mechanism is physically distinct
from other quantum gravity proposals:

In \emph{loop quantum gravity}~\cite{rovelli_2004},
the singularity is resolved by quantum geometry:
area and volume operators have discrete spectra with
minimum eigenvalues of order $\lP^2$ and $\lP^3$.

In \emph{string theory}~\cite{polchinski_1998}, the
finite size of the string ($\ell_s \sim \lP$) prevents
localization below the string scale.

In the \emph{superfluid vacuum}, the singularity is
resolved by the Bohm potential of the condensate:
the quantum pressure of the vacuum fluid prevents
gravitational collapse below the healing length. The
mechanism is hydrodynamic, not geometric or
string-theoretic, but produces the same minimum
length $\xi \sim \lP$.

The key distinction is that $\lP$ is \emph{derived}
from the condensate parameters, not introduced as a
fundamental scale. The ``quantization of geometry''
is an emergent consequence of the quantum coherence
of the vacuum condensate.


% ====================================================================
\section{Unification of Fundamental Constants}
\label{sec:unification}
% ====================================================================

The identification $\xi = \lP$ provides a unified
interpretation of the three fundamental constants
$\hbar$, $c$, and $G$ as properties of a single
physical substrate---the vacuum condensate:

\begin{itemize}
\item $c = c_s$: the speed of light is the sound speed
  of the condensate at the cosmological background
  density~\cite{kulangiev_2026a}. It is determined by
  the logarithmic equation of state.

\item $G \sim c^2/(\xi^2\rho_0)$: Newton's constant
  is the emergent gravitational coupling, fixed by
  the healing length $\xi = \lP$ and the vacuum mass
  density $\rho_0 \sim \mP/\lP^3$ through the
  dimensional structure of the linearized condensate
  response (Sec.~\ref{sec:derivation}).

\item $\hbar = \heff$: Planck's constant governs the
  quantum coherence of the condensate. It sets the
  healing length $\xi = \hbar/(\meff c)$ and the
  Gausson width. The de Broglie relation~\cite{de_broglie_1927}
  $p = \hbar k$ is the dispersion relation of
  condensate excitations.
\end{itemize}

The Planck units are then:
\begin{align}
    \lP &= \xi = \text{healing length of the vacuum}\,,
    \label{eq:lP}\\
    \tP &= \xi/c = \text{phase crossing time}\,,
    \label{eq:tP}\\
    \mP &\sim \meff = \text{constituent mass of the vacuum}\,,
    \label{eq:mP}\\
    E_{\mathrm{P}} &= \mP c^2 = \text{rest energy of one constituent}\,.
    \label{eq:EP}
\end{align}
The Planck scale is not a fundamental discreteness
of spacetime but the scale at which the emergent
spacetime description (acoustic metric) gives way
to the underlying quantum dynamics of the vacuum
condensate.


% ====================================================================
\section{Discussion}
\label{sec:discussion}
% ====================================================================

\subsection{What is derived and what is assumed}

The derivation of $\xi = \lP$ rests on four inputs:

(i) The dimensional form of the gravitational
coupling, $G \sim c^2/(\xi^2\rho_0)$
(Eq.~\ref{eq:G}), which follows from dimensional
analysis once the available parameters of the
emergent-gravity sector are fixed
($c$, $\rho_0$, and a length scale).

(ii) The healing-length structure of the soliton,
$\xi = \heff/(\meff c)$ (Eq.~\ref{eq:healing}),
which arises from balancing the quantum potential
against the logarithmic interaction at the soliton
core.

(iii) The natural packing density $n_0 \sim 1/\xi^3$
(Eq.~\ref{eq:n0}), reflecting one condensate
particle per healing volume.

(iv) The identification $\heff = \hbar$, the
simplest assumption consistent with the condensate's
excitations being identified with standard quantum
particles.

Of these, (i) and (ii) follow from dimensional
analysis once the framework's structure is fixed,
while (iii) and (iv) are physical assumptions. The
result $\meff \sim \mP$, $\xi \sim \lP$,
$\rho_0 \sim \mP/\lP^3$ is robust to $O(1)$ factors
throughout the chain.

\subsection{The vacuum density}

From the self-consistent
identifications~(\ref{eq:planck_identifications}),
the condensate has one particle of mass $\meff \sim \mP$ per Planck volume, giving a mass density at the
Planck scale:
\begin{equation}
    \rho_0 \;=\; \meff\,n_0 \;\sim\; \frac{\mP}{\lP^3}
    \;\approx\; 5.2 \times 10^{96}\;\mathrm{kg/m^3}\,.
\label{eq:rho0}
\end{equation}
This is the Planck density---the natural maximum
density for any relativistic quantum field. It is,
however, the density of the \emph{vacuum condensate}
itself, not of ordinary matter. A useful intuition:
the condensate has bulk modulus $K \sim \rho_0 c^2
\sim 10^{113}\;\mathrm{Pa}$, larger than that of
diamond by $\sim 10^{101}$ orders of magnitude.
Gravity is weak because the vacuum is stiff: a
massive source produces only a tiny fractional
density perturbation in the condensate, and this
small perturbation is what drives gravitational
acceleration. The energy density of the background
itself, $\rho_0 c^2 \sim 10^{113}\;\mathrm{J/m^3}$,
is comparable to the QFT zero-point estimate of
the cosmological constant; in the present framework
this is not directly observable, since gravitational
effects are sourced by perturbations
$\delta\rho$ rather than by the background density
$\rho_0$.

\newpage
\subsection{Testable consequences}

The healing-length regularization makes two
predictions that distinguish it from classical GR:

(i) \emph{No singularity at $r = 0$}: The density
inside a black hole is bounded by
$\rho_{\max} \sim e\,\rho_0$. This affects the
information paradox: if no singularity forms, the
interior is a finite-density condensate core, and
unitarity may be preserved.

(ii) \emph{Modified dispersion at high energies}:
At energies approaching $E_{\mathrm{P}} = \mP c^2$,
the phonon dispersion relation deviates from
$\omega = c k$ due to the Bohm potential
corrections. This is analogous to the
roton-maxon spectrum in superfluid
helium~\cite{zloshchastiev_2012} and is in
principle testable through gamma-ray burst
time-of-flight measurements~\cite{amelino_2013}.

\subsection{Limitations}

The derivation assumes a single species of condensate
particle with mass $\meff \sim \mP$. If the condensate
has a more complex structure (e.g., multiple species
or a composite fundamental unit), the healing length
could differ from $\lP$ by numerical factors. The $O(1)$ uncertainty in the Gausson energy integral
propagates to an $O(1)$ factor in $\xi/\lP$,
which is within the expected precision of the
order-of-magnitude analysis.

The interior structure of black holes in the
superfluid framework requires solving the full
nonlinear wave equation~(\ref{eq:hj}) with $Q$
retained, which is beyond the scope of this paper.
The regularization argument establishes the existence
of a maximum density and minimum length, but not the
detailed density profile inside the horizon.


% ====================================================================
\section{Conclusion}
\label{sec:conclusion}

We have shown that the logarithmic superfluid vacuum
framework provides a complete UV regularization of
the two singular limits of general relativity:

(1) \emph{Spatial singularity.} The quantum potential
$Q = \Box\sqrt{\rho}/\sqrt{\rho}$ prevents
gravitational collapse below the healing length
$\xi \sim \lP$. The density is bounded by
$\rho_{\max} \sim e\,\rho_0$, and the Kretschmann
scalar is bounded by $K \lesssim 1/\lP^4$.

(2) \emph{Temporal singularity.} The chemical
potential $\mu(x) = \meff c_s^2/\heff$, which sets
the local clock rate, is bounded from both sides
by the bounded density. The clock never stops
($\mu > 0$), time dilation is always finite, and
the minimum resolvable time interval is
$\Delta t_{\min} \sim \tP$.

Both regularizations arise from the same physical
mechanism---the Bohm potential of the vacuum
condensate---and produce the same characteristic
scale---the Planck scale---which emerges as the
healing length $\xi = \hbar/(\meff c) \sim \lP$,
not as an assumed discreteness.

The framework unifies $\hbar$, $c$, and $G$ as
properties of a single substrate: the speed of
sound, the quantum coherence scale, and the
emergent gravitational coupling of the vacuum
condensate. The Planck mass $\mP \sim \meff$ is
the constituent mass, and the Planck length is the
coherence length at which the emergent spacetime
gives way to the underlying quantum dynamics.


% ====================================================================
\section*{Acknowledgments}
% ====================================================================

Claude (Anthropic) and Gemini (Google) were used as editorial and computational tools during manuscript preparation. All scientific content, equations, and interpretations are the author's own.


\begin{thebibliography}{25}

\bibitem{kretschmann_1917}
E.~Kretschmann,
\"Uber den physikalischen Sinn der Relativit\"{a}tspostulate, A.\ Einsteins neue und seine urspr\"{u}ngliche Relativit\"{a}tstheorie,
Ann.\ Phys.\ (Leipzig) \textbf{358}, 575 (1917).

\bibitem{madelung_1927}
E.~Madelung,
Quantentheorie in hydrodynamischer Form,
Z.\ Phys.\ \textbf{40}, 322 (1927).

\bibitem{bohm_1952}
D.~Bohm,
A suggested interpretation of the quantum theory in terms of ``hidden'' variables. I,
Phys.\ Rev.\ \textbf{85}, 166 (1952).

\bibitem{de_broglie_1927}
L.~de~Broglie,
La m\'{e}canique ondulatoire et la structure atomique de la mati\`{e}re et du rayonnement,
J.\ Phys.\ Radium \textbf{8}, 225 (1927).

\bibitem{ashtekar_2006}
A.~Ashtekar, T.~Pawlowski, and P.~Singh,
Quantum nature of the big bang: Improved dynamics,
Phys.\ Rev.\ D \textbf{74}, 084003 (2006).

\bibitem{garay_1995}
L.~J.~Garay,
Quantum gravity and minimum length,
Int.\ J.\ Mod.\ Phys.\ A \textbf{10}, 145 (1995).

\bibitem{penrose_1965}
R.~Penrose,
Gravitational collapse and space-time singularities,
Phys.\ Rev.\ Lett.\ \textbf{14}, 57 (1965).

\bibitem{hawking_1970}
S.~W.~Hawking and R.~Penrose,
The singularities of gravitational collapse and cosmology,
Proc.\ Roy.\ Soc.\ Lond.\ A \textbf{314}, 529 (1970).

\bibitem{dewitt_1967}
B.~S.~DeWitt,
Quantum theory of gravity. I. The canonical theory,
Phys.\ Rev.\ \textbf{160}, 1113 (1967).

\bibitem{rovelli_2004}
C.~Rovelli,
\emph{Quantum Gravity} (Cambridge, 2004).

\bibitem{polchinski_1998}
J.~Polchinski,
\emph{String Theory} (Cambridge, 1998).

\bibitem{kulangiev_2026a}
B.~B.~Kulangiev,
Conditions for emergent gravitational light bending
from a logarithmic superfluid vacuum,
Preprint (2026).
\url{https://kulangiev.github.io#paper1}

\bibitem{volovik_2003}
G.~E.~Volovik,
\emph{The Universe in a Helium Droplet} (Oxford, 2003).

\bibitem{bialynicki_birula_1979}
I.~Bia{\l}ynicki-Birula and J.~Mycielski,
Gaussons: Solitons of the logarithmic Schr\"{o}dinger
equation,
Phys.\ Scr.\ \textbf{20}, 539 (1979).

\bibitem{zloshchastiev_2012}
K.~G.~Zloshchastiev,
Volume element structure and roton-maxon-phonon excitations in superfluid helium beyond the Gross-Pitaevskii approximation,
Eur.\ Phys.\ J.\ B \textbf{85}, 273 (2012).

\bibitem{amelino_2013}
G.~Amelino-Camelia,
Quantum-spacetime phenomenology,
Living Rev.\ Rel.\ \textbf{16}, 5 (2013).

\end{thebibliography}

\end{document}
